% AUTO-GENERATED FILE. Source: pub/claims.toml
Evaluation of post-hoc explanations in machine learning remains fragmented
across incompatible metric families, yielding method comparisons that are
sensitive to which endpoints are chosen. We address this through two
complementary contributions. First, drawing on a 44-paper structured
scoping corpus assembled through a documented five-database search and
transparent screening record, we build a four-axis taxonomy---by
evaluation target, evidence source, quality property, and task
context---that maps the measurement landscape and exposes three recurring
gaps: proxy metrics dominate despite not
capturing semantics; human-grounded constructs remain underspecified; and
semantic evaluation is growing faster than its empirical validation base.
Second, we fill the proxy-layer gap with a paired empirical benchmark
comparing LIME and SHAP across 75 matched cells spanning five model
families, five seeds, and three sampling sizes on the Adult tabular
benchmark. Within this protocol, results are consistent and large in
magnitude: SHAP dominates all measured quality endpoints (stability
$d_z{=}3.00$, fidelity $d_z{=}4.82$, faithfulness gap $d_z{=}2.63$),
while LIME is faster for non-tree model families (median 53.3 ms vs.694.6 ms for SHAP); for tree-based models, SHAP's 	exttt{TreeExplainer}
reverses this ordering. A tabular cross-dataset
extension on two additional datasets provides partial support for the
fidelity direction, but does not establish cross-modal generality.
Together, the taxonomy and the benchmark motivate a
model-architecture-conditioned deployment pattern for tabular
classification: SHAP is preferred at both fidelity and latency for
tree-based models, while LIME retains a latency advantage for black-box
architectures lacking a model-aware explainer. All experimental artifacts
are publicly available.
